\section{Metodología}

\subsection{Enfoque propuesto}

\textbf{Formalización del problema.} Siguiendo a Mitchell \cite{mitchell1997machine}, se define la terna $\langle T, E, P \rangle$. La \textbf{tarea} $T$ es inducir $f:\mathcal{X}\rightarrow\{0,1\}$ que asigne a cada hogar $i$ la etiqueta $y_i = \mathbb{I}(\texttt{POBREZA}_i \in \{1,2\})$, es decir, condición de pobreza extrema (código 1) o no extrema (código 2), agrupadas como pobreza monetaria total frente a no pobre (código 3), según el Módulo 34 (Sumaria) del INEI. La \textbf{experiencia} $E$ es la ENAHO 2024 de \texttt{DOMINIO = 8}: $N = 4{,}090$ hogares, de los cuales 774 son pobres (18.92\%), con razón de desbalance $\gamma \approx 4.28$. La prueba fuera de tiempo es la ENAHO 2025 sin los hogares panel: 3,199 hogares (17.4\% pobres). La \textbf{medida de desempeño} $P$ incluye la PR-AUC, el error de exclusión $\mathrm{EE} = FN/(TP+FN)$ y el de inclusión $\mathrm{IE} = FP/(TP+FP)$ con cobertura fija $B$, la exhaustividad, la precisión y el puntaje de Brier tras recalibrar las probabilidades. Los clasificadores minimizan una pérdida logística sensible al costo:
\begin{equation}
\mathcal{L}(\theta) = -\frac{1}{N}\sum_{i=1}^{N}\bigl[\alpha\, y_i \ln \hat{p}_i + (1-y_i)\ln(1-\hat{p}_i)\bigr] + \lambda\,\Omega(\theta),
\end{equation}
donde $\hat{p}_i = P(y_i = 1 \mid \mathbf{x}_i;\theta)$, $\Omega$ es la regularización y $\alpha \in \{1, 2, \gamma\}$ es un hiperparámetro elegido por validación ($\gamma$ es un cociente de frecuencias, no un costo social). Como reponderar equivale a desplazar el umbral de un modelo sin costos a $p^* = 1/(1+\alpha)$ \cite{elkan2001foundations}, no se aplican ambas correcciones a la vez.

\textbf{Comportamiento entrada/salida.} La entrada es $\mathbf{x}_i = [\mathbf{x}_i^{(\text{viv})}, \mathbf{x}_i^{(\text{dem})}, \mathbf{x}_i^{(\text{edu})}, \mathbf{x}_i^{(\text{emp})}]$, con 28 variables observables:
\begin{itemize}
    \item Vivienda (Módulo 01): materiales, habitaciones, tenencia, agua, desagüe, combustible e internet.
    \item Demografía (Módulo 02): tamaño, menores, adultos mayores y dependencia; sexo y edad del jefe.
    \item Educación (Módulo 03): nivel del jefe y máximo del hogar.
    \item Empleo (Módulo 05): ocupados, categoría ocupacional, horas y afiliación a pensión del jefe.
\end{itemize}
Ninguna variable de ingreso o gasto entra en $\mathbf{x}_i$; el gasto solo es la variable dependiente de la línea base. La salida es un puntaje por hogar ($\hat{p}_i$, o el gasto predicho con signo invertido en la línea base). La decisión selecciona al 20\% de hogares con mayor puntaje, que es la comparación principal, o usa un umbral $\tau$ fijado en la validación de 2024 y congelado para 2025 (H3).

\textbf{Operadores y adaptaciones al problema.} La Figura~\ref{fig:pipeline} muestra el flujo y la Tabla~\ref{tab:adaptaciones} el sustento en la literatura. Las adaptaciones 1--3 son decisiones de ingeniería de datos motivadas por la jerarquía censal de la ENAHO (vivienda $\rightarrow$ hogar $\rightarrow$ individuo; la agregación $\Phi$ sigue las recomendaciones de diseño de registros de Grosh et al. \cite{grosh2022revisiting}), mientras que las adaptaciones 4--8 son formulaciones algorítmicas sustentadas en la literatura:
\begin{enumerate}
    \item \textit{Imputación intra-vivienda.} En 2024, 72 hogares secundarios (1.76\%) que comparten vivienda con el hogar principal no tienen datos físicos (paredes, pisos, habitaciones). Los heredan de su vivienda $(\texttt{CONGLOME}, \texttt{VIVIENDA})$ mediante relleno hacia adelante y atrás restringido a cada grupo, sin imputar valores globales.
    \item \textit{Agregación individuo$\rightarrow$hogar.} Un operador $\Phi$ toma las personas $j$ del hogar $i$ y produce dos ramas: los atributos del jefe ($P203 = 1$) y funcionales del hogar, por ejemplo $\text{dep}_i = \sum_j \mathbb{I}(e_j<15 \lor e_j \ge 65) / \max\bigl(1, \sum_j \mathbb{I}(15 \le e_j < 65)\bigr)$, donde $e_j$ es la edad de la persona $j$. Se agregan el máximo nivel educativo, la proporción de ocupados y las personas por habitación.
    \item \textit{Agrupación de categorías raras ajustada solo en entrenamiento.} Las categorías con menos de 1\% se fusionan por afinidad. En pisos (2024): acabado noble 49.8\% de los hogares (pobreza 8.2\%), cemento 46.0\% (28.9\%) y precario 4.2\% (37.2\%). Las tasas se calculan solo con 2024 para evitar fuga de la etiqueta.
    \item \textit{Línea base econométrica.} PMT por mínimos cuadrados: $\ln(\text{gasto}_i) = \mathbf{w}^{\top}\mathbf{x}_i + \varepsilon_i$, con $\hat{y}_i = \mathbb{I}\bigl(\widehat{\ln \text{gasto}_i} < \ln \text{LINEA}\bigr)$. Se compara con una regresión logística ElasticNet, un \textit{random forest} y LightGBM \cite{ke2017lightgbm}, cuyos hiperparámetros se eligen por validación agrupada.
    \item \textit{Comparación a igual cobertura.} Todos los modelos seleccionan la misma proporción $B = 20$\% de hogares, y se trazan curvas exclusión-inclusión para $B$ entre 10\% y 40\%. Así se aísla el aporte del modelo de la elección del umbral.
    \item \textit{Protocolo contra la fuga de información.} (a) Se excluyen del test los 930 hogares de 2025 que también fueron encuestados en 2024. (b) La validación interna usa \texttt{GroupKFold} de 5 particiones por conglomerado, porque 702 de los 704 conglomerados se repiten entre años. (c) Codificación, escalamiento y agrupación de categorías se ajustan dentro de un \texttt{Pipeline} solo con datos de entrenamiento.
    \item \textit{Inferencia.} Las diferencias entre modelos (H1, H2, H4) se estiman con intervalos de confianza por \textit{bootstrap} de conglomerados (1,000 réplicas), en lugar de la prueba de McNemar, que no compara PR-AUC ni F1.
    \item \textit{Interpretabilidad.} TreeSHAP \cite{lundberg2020local} descompone el puntaje de cada hogar en contribuciones aditivas, $f(\mathbf{x}_i) = \phi_0 + \sum_{j=1}^{d}\phi_j(\mathbf{x}_i)$. Así se verifica que las variables más influyentes sean coherentes con la literatura y se detecta cualquier variable que filtre información del gasto.
\end{enumerate}
Un chequeo preliminar con las 28 variables descritas, cuyo script está en el repositorio, dio $\mathrm{EE} \approx 0.41$ con $B = 20$\% para la línea base, la regresión logística y el \textit{random forest}. Esto motiva formular H2 como una prueba de equivalencia. Las métricas se reportan por hogar y, como robustez, ponderadas por personas (\texttt{FACTOR07} $\times$ miembros).

\begin{table}[htbp]
\centering
\caption{Adaptaciones del enfoque y publicación que las sustenta.}
\label{tab:adaptaciones}
\fontsize{6.8}{8.2}\selectfont
\setlength{\tabcolsep}{3pt}
\begin{tabular}{L{2.5cm}L{3.1cm}L{1.4cm}}
\toprule
\textbf{Adaptación} & \textbf{Cómo se aplica} & \textbf{Sustento} \\
\midrule
Agregación del hogar ($\Phi$) & Resumen de jefatura y carga & \cite{grosh2022revisiting} \\
\addlinespace[2pt]
Línea base PMT-MCO & $\ln(\text{gasto})$ con corte en la línea & \cite{brown2018poor,mcbride2018retooling} \\
\addlinespace[2pt]
Pérdida sensible al costo & $\alpha$ como hiperparámetro; sin doble corrección del umbral & \cite{elkan2001foundations} \\
\addlinespace[2pt]
Igual cobertura y curva EE-IE & Selección de 20\% de hogares & \cite{noriega2020algorithmic} \\
\addlinespace[2pt]
Prueba fuera de tiempo sin panel & Entrenar 2024, evaluar 2025 sin 930 hogares & \cite{aiken2023moving} \\
\addlinespace[2pt]
Validación agrupada & \texttt{GroupKFold} por conglomerado & \cite{roberts2017cross} \\
\addlinespace[2pt]
Ensambles de árboles & \textit{Random forest} y LightGBM & \cite{grinsztajn2022tree,ke2017lightgbm} \\
\addlinespace[2pt]
Interpretabilidad & TreeSHAP por hogar & \cite{lundberg2020local} \\
\bottomrule
\end{tabular}
\end{table}

\begin{figure}[htbp]
\centering
\begin{tikzpicture}[
    basebox/.style={rectangle, rounded corners=2.5pt, line width=0.7pt, inner sep=3pt,
        font=\fontsize{6.6}{7.9}\selectfont, align=left, text width=0.92\columnwidth},
    boxgray/.style={basebox, draw=gray!70, fill=black!4},
    boxgreen/.style={basebox, draw=primarygreen, fill=primarygreen!7},
    boxteal/.style={basebox, draw=secondaryteal, fill=secondaryteal!7},
    arrow/.style={-{Stealth[length=4pt, width=3pt]}, line width=0.8pt, draw=primarygreen}
]
\node[boxgray] (ingesta) {
    \textbf{\color{primarygreen}1. Datos: ENAHO 2024--2025, \texttt{DOMINIO = 8}}\\
    $\bullet$ Hogar: Módulo 01 (vivienda) y Módulo 34 (etiqueta $y_i$).\\
    $\bullet$ Persona: Módulos 02, 03 y 05 (demografía, educación, empleo).
};
\node[boxgreen, below=4pt of ingesta] (preproc) {
    \textbf{\color{primarygreen}2. Representación del hogar (sin ingreso ni gasto)}\\
    $\bullet$ Imputación intra-vivienda; agregación $\Phi$ (jefe y hogar).\\
    $\bullet$ Agrupación de categorías raras ajustada solo con 2024.
};
\node[boxteal, below=4pt of preproc] (modelado) {
    \textbf{\color{secondaryteal}3. Modelos}\\
    $\bullet$ Línea base PMT-MCO sobre $\ln(\text{gasto})$; logística ElasticNet.\\
    $\bullet$ \textit{Random forest} y LightGBM con pérdida sensible al costo ($\alpha$).
};
\node[boxgray, below=4pt of modelado] (validacion) {
    \textbf{\color{primarygreen}4. Validación y evaluación}\\
    $\bullet$ 2024: \texttt{GroupKFold} por conglomerado (hiperparámetros, $\alpha$, $\tau$).\\
    $\bullet$ 2025 sin 930 hogares panel: PR-AUC, EE e IE con $B = 20$\%,\\
    \hphantom{$\bullet$ }curvas exclusión-inclusión e IC por \textit{bootstrap}.
};
\draw[arrow] (ingesta) -- (preproc);
\draw[arrow] (preproc) -- (modelado);
\draw[arrow] (modelado) -- (validacion);
\end{tikzpicture}
\caption{Flujo metodológico. $y_i$: pobreza oficial del hogar; $\Phi$: agregación de personas a hogar; $\alpha$: peso de la clase pobre en la pérdida; $\tau$: umbral de decisión; $B$: proporción de hogares seleccionados; EE e IE: errores de exclusión e inclusión.}
\label{fig:pipeline}
\end{figure}
